% !TeX spellcheck = pt_BR

\chapter*{Do que se trata a Cinemática}
\addcontentsline{toc}{chapter}{Do que se trata a Cinemática}



%Escrever aqui uma boa intro sobre o que se trata a cinemática, e que tal estudo aqui ficará dividido em duas partes: a parte unidimensional, que é escalar, e a parte bidimensional, que é vetorial.





A Cinemática é a área da física que trata do estudo do movimento. Buscamos aqui compreender as relações entre movimentos de um corpo em relação a outro, da rapidez com que um objeto se move e do quanto sua velocidade aumenta ou diminui.

Por ora vamos analisar o movimento dos corpos de forma pura, sem se preocupar com o que causa tal movimento. Historicamente, foi pela Cinemática que a Física começou a se formar, com os estudos de Galileu sobre a queda dos corpos e o movimento dos projéteis, buscando entender como tais movimentos poderiam ser descritos, sem se preocupar com o que causava tais movimentos -- as forças.

Ao suspender, provisoriamente, a discussão sobre forças, a Cinemática permite construir uma linguagem precisa para falar de movimento. Essa separação não é artificial: antes de explicar o mundo, é preciso aprender a descrevê-lo, e é isso que faremos por meio deste estudo.

O nosso estudo se dividirá em duas partes: a Cinemática Escalar, que tratará essencialmente dos tipos mais simples de movimento: os movimentos unidimensionais. Neste contexto, como teremos uma única dimensão na qual o movimento pode ocorrer, o caráter vetorial de grandezas como velocidade e aceleração poderá ficar de lado, haja vista que, para um corpo que se move horizontalmente, restam-nos duas opções para onde ele pode estar se movendo: para a direita ou para a esquerda. Essas duas opções podem ser distinguidas por meio de uma atribuição de sinais: uma velocidade positiva pode representar um corpo se movendo para direita, enquanto uma velocidade negativa representaria um corpo que se move para esquerda.

Dessa forma, usando alguns artifícios algébricos como este, poderemos estudar em detalhe diversos tipos de movimentos em uma dimensão sem necessariamente precisar fazer um tratamento vetorial das grandezas cinemáticas envolvidas -- essa é a Cinemática Escalar.

Posteriormente, quando quisermos estudar movimentos bidimensionais ou mesmo tridimensionais, seremos obrigados a usar o tratamento vetorial: existem apenas dois tipos de sinais -- positivo ou negativo -- porém a velocidade agora pode apontar para uma infinidade de direções distintas, que não poderão mais ser diferenciadas apenas com a atribuição de um sinal algébrico. Com isso, adentraremos no estudo da Cinemática Vetorial, em que redefiniremos todas as grandezas outrora estudadas na Cinemática Escalar, porém agora dando o devido tratamento vetorial a elas. Essa parte começará com uma breve apresentação sobre os vetores, suas características geométricas e suas operações. Com isso, estaremos munidos de todas as ferramentas necessárias para descrever a evolução de qualquer tipo de movimento -- e teremos, em mãos, tudo que precisamos para avançar nos estudos da Dinâmica, que constitui a segunda parte do nosso estudo.







